\documentclass[../DoAn.tex]{subfiles}
\begin{document}

\begin{center}
    \Large{\textbf{ABSTRACT}}
\end{center}
\vspace{1cm}

The research project ``Building a Multi-Platform Automated Grading System using OMR (Optical Mark Recognition)'' investigates and develops a comprehensive system to automate the multiple-choice exam grading process, thereby significantly reducing time consumption and errors compared to traditional manual grading. In the context of increasingly widespread adoption of multiple-choice examinations at high schools and universities in Vietnam with thousands of candidates per session, manual grading is not only labor-intensive but also prone to fatigue-induced errors. Existing OMR solutions on the market predominantly require high licensing costs, rely on physical scanning hardware, or offer command-line interfaces that are difficult to operate and lack mobile platform support.

The proposed system comprises three principal components developed on modern technology platforms. The first component is a mobile application built with Flutter and Dart, enabling users to capture images of answer sheets using smartphone cameras and recognize answers through OMR technology. The distinctive feature of this component is that the OMR recognition engine is deployed entirely on the mobile device (client-side), allowing immediate grading without internet connectivity, while supporting an offline-first architecture with automatic synchronization upon network availability. The second component is a web application built with React and TypeScript, providing a management interface for teachers to create exams, manage question banks, monitor grading results, and export statistical reports. The third component is a backend developed with Node.js and Express, combined with MongoDB database, providing RESTful APIs for the entire system and integrating with AMC (Auto Multiple Choice) to automatically generate multiple exam versions along with precisely calibrated OMR coordinates.

The proposed solution achieves several notable contributions. Firstly, the system integrates an automatic exam generation pipeline using AMC, enabling the generation of N distinct versions with fully shuffled questions and answers, while simultaneously exporting a JSON template containing calibrated bubble coordinates, ensuring absolute accuracy without reliance on manual calculations. Secondly, the mobile-based OMR recognition engine achieves an average accuracy exceeding 95\% on camera-captured images under normal lighting conditions, with an average processing time below two seconds per sheet on mid-range devices. Thirdly, the system incorporates an artificial intelligence module using Google Gemini API to analyze examination results, identify topics where students demonstrate weaknesses, and recommend suitable practice questions. Fourthly, thanks to its offline-first architecture, the system operates reliably in exam room environments with limited or no network connectivity, ensuring zero data loss.

Experimental results on a dataset comprising over one thousand exam sheets from actual testing sessions demonstrate that the system achieves approximately 96.5\% recognition accuracy on scanned answer sheets, with an average processing time of 1.8 seconds per sheet on mid-range devices. The system has been deployed in a pilot program involving over one hundred and fifty users and has received positive feedback from both teachers and students regarding its usability and practical effectiveness.

\textbf{Keywords:} OMR, automated grading, Flutter, React, Node.js, MongoDB, offline-first, image recognition

\end{document}
